\subsection{Rationale Re-Kaskade am Kepler-Anker}\label{sec:eabc-zeugen-rational}

Am Fixpunkt $M=113160$ liefern die kanonischen rationalen EABC-Zeugen ($a=0$, $M/\mathrm{Norm}\in\mathbb{Z}$) eine absteigende Schalenkaskade über die Normen $1,\ldots,6$. Die Impuls-Eichung
\[
  \Lambda_{\mathrm{impuls}} = \frac{Re_{\mathrm{krit}}}{M/\mathrm{Norm}_{\mathrm{Ref}}},
  \qquad
  Re = \Lambda_{\mathrm{impuls}}\cdot\frac{M}{\mathrm{Norm}},
\]
wird an der tiefsten verdichteten Schale $\mathrm{Norm}=6$ gegen den klassischen laminar-turbulenten Umschlag $Re_{\mathrm{krit}}\approx 2300$ kalibriert. Es ergibt sich $\Lambda_{\mathrm{impuls}}\approx 0{,}12195122$ (vgl.\ \texttt{Reynolds.py}).

\begin{table}[ht]
  \centering
  \caption{Kanonische Fixpunkte des rationalen Sektors ($a=0$) und $Re$-Eichung mit $\Lambda_{\mathrm{impuls}}\approx 0{,}12195122$.}
  \label{tab:eabc-rational}
  \begin{tabular}{rrrrr}
    \toprule
    $(c_e,c_a,c_b,c_c)$ & $\mathrm{Norm}$ & $M/\mathrm{Norm}$ & $Re$ & Physikalischer Status \\
    \midrule
    $(1, 0, 0, 0)$ & 1.0 & 113160 & 13800.00 & Laminarer Fixpunkt \\
    $(2, 0, 0, 0)$ & 2.0 & 56580 & 6900.00 & Laminarer Fixpunkt \\
    $(2, 0, 2, 1)$ & 3.0 & 37720 & 4600.00 & Laminarer Fixpunkt \\
    $(4, 0, 0, 0)$ & 4.0 & 28290 & 3450.00 & Laminarer Fixpunkt \\
    $(4, 0, 3, 0)$ & 5.0 & 22632 & 2760.00 & Laminarer Fixpunkt \\
    $(4, 0, 4, 2)$ & 6.0 & 18860 & 2300.00 & Kritische Grenze (Umschlag) \\
    \bottomrule
  \end{tabular}
\end{table}

Die höheren harmonischen Stufen $Re\in\{4600,6900,13800\}$ markieren im Forschungsprogramm stabil-laminare Resonanzbauten oberhalb der viskosen Grenze; die Schale $\mathrm{Norm}=6$ fixiert den Umschlag exakt auf $Re_{\mathrm{krit}}$.
